% ---------------------------------------------------------------------------
% Theorem-like environments
% ---------------------------------------------------------------------------
\newtheorem{theorem}{\sc\bf Theorem}
\newtheorem{example}{\sc\bf Example}
\newtheorem{definition}[theorem]{\sc\bf Definition}

\newcommand{\proof}{\noindent{\sc\bf Proof}\quad }
\newcommand{\bbox}{\phantom{1}\hfill{\rule{6pt}{6pt}}}
\newcommand{\remark}[1]{\noindent{\sc Remark}:\ #1}
\newcommand{\hint}[1]{{\sc Hint}:\ {\it #1}}
\newcommand{\note}[1]{({\sc Note}:\ #1)}

% ---------------------------------------------------------------------------
% Headings
% ---------------------------------------------------------------------------

% Unnumbered subheading, listed in the PDF bookmarks but not in the printed TOC.
\newcommand{\boxit}[1]{%
  \phantomsection
  \addcontentsline{toc}{subsubsection}{#1}%
  \subsection*{#1}}

% "Exercises" subheading at the end of a section.
\newcommand{\exercises}{%
  \phantomsection
  \pdfbookmark[2]{Exercises}{exercises.\thesection}%
  \label{exerstart:\thesection}%
  \subsection*{Exercises}%
  \noindent}

\newcommand{\technology}{\bigskip\noindent\fbox{\color{red}USING
TECHNOLOGY}\par}

% ---------------------------------------------------------------------------
% Lists
% ---------------------------------------------------------------------------
\newlist{exerciselist}{enumerate}{1}
\setlist[exerciselist]{label=\textbf{\arabic*.},ref=\arabic*}
\newlist{alist}{enumerate}{1}
\setlist[alist]{label=\textbf{(\alph*)},ref=\alph*}
\newlist{rmlist}{enumerate}{1}
\setlist[rmlist]{label=\textbf{(\roman*)},ref=\roman*}
\newlist{numberlist}{enumerate}{1}
\setlist[numberlist]{label=\textbf{\arabic*.},ref=\arabic*}

% Interrupt an exercise list with a line of text; numbering continues.
\newcommand{\exercisetext}[1]{\end{exerciselist}{\itshape #1}\begin{exerciselist}[resume]}

% ---------------------------------------------------------------------------
% Exercises and solutions
% ---------------------------------------------------------------------------
\renewcommand{\part}[1]{{\bf (#1)}}
\newcommand{\answer}[1]{\;{\bf $\mathbf{#1}$ (p.~\pageref{exer:#1})}\;}
\newcommand{\solution}{\vskip10pt\noindent{\sc\bf Solution }\  }
\newcommand{\solutionpart}[1]{\vskip10pt\noindent\underbar{\color{blue}\sc
Solution({\bf #1})\ }}

% Markers for exercises that call for computation and graphics (C, C/G)
% or for extensive use of technology (L).
\newcommand{\Cex}{\fbox{\textcolor{red}{C}}\, }
\newcommand{\CGex}{\fbox{\textcolor{red}{C/G}}\, }
\newcommand{\Lex}{\fbox{\textcolor{red}{L}}\, }

% ---------------------------------------------------------------------------
% Math shorthands
% ---------------------------------------------------------------------------
\newcommand{\dst}{\displaystyle}
\newcommand{\lims}[2]{\,\bigg|_{#1}^{#2}}

% Column vectors and matrices
\newcommand{\twocol}[2]{\left[\begin{array}{l}#1\\#2\end{array}\right]}
\newcommand{\ctwocol}[2]{\left[\begin{array}{c}#1\\#2\end{array}\right]}
\newcommand{\threecol}[3]{\left[\begin{array}{r}#1\\#2\\#3\end{array}\right]}
\newcommand{\cthreecol}[3]{\left[\begin{array}{c}#1\\#2\\#3\end{array}\right]}
\newcommand{\col}[2]{\left[\begin{array}{c}#1_1\\#1_2\\\vdots\\#1_#2\end{array}\right]}
\newcommand{\colfunc}[2]{\left[\begin{array}{c}#1_1(t)\\#1_2(t)\\\vdots\\#1_#2(t)\end{array}\right]}
\newcommand{\twobytwo}[4]{\left[\begin{array}{rr}
#1&#2\\#3&#4\end{array}\right]}
\newcommand{\threebythree}[9]{\left[\begin{array}{rrr}
#1&#2&#3\\#4&#5&#6\\#7&#8&#9\end{array}\right]}
\newcommand{\cthreebythree}[9]{\left[\begin{array}{ccc}
#1&#2&#3\\#4&#5&#6\\#7&#8&#9\end{array}\right]}
\newcommand{\matfunc}[3]{\left[\begin{array}{cccc}#1_{11}(t)&#1_{12}(t)&\cdots
&#1_{1#3}(t)\\#1_{21}(t)&#1_{22}(t)&\cdots&#1_{2#3}(t)\\\vdots&
\vdots&\ddots&\vdots\\#1_{#21}(t)&#1_{#22}(t)&\cdots&#1_{#2#3}(t)
\end{array}\right]}

% Characteristic determinant
\newcommand{\twochar}[4]{\left|\begin{array}{cc}
#1-\lambda&#2\\#3&#4-\lambda\end{array}\right|}

% ---------------------------------------------------------------------------
% Layout helpers
% ---------------------------------------------------------------------------
\newcommand{\bigbox}[1]{\vskip10pt\vbox{\hrule\hbox{\vrule\kern5pt
\vbox{\noindent
\strut\kern5pt\mbox{}\newline#1\mbox{}\newline
\kern5pt\strut}\kern5pt\vrule}\hrule}}
\newcommand{\blankbox}[1]{\vskip10pt\vbox{\hbox{\kern5pt
\vbox{\noindent\strut\kern5pt\vspace*{-15pt}#1\vspace*{-20pt}
\kern5pt\strut}\kern5pt}}}

\newenvironment{newtable}{\begin{quote}
Table\space\refstepcounter{table}\thetable.}{\end{quote}}

% Text row inside an array; must come right after a line break (\\).
\newcommand{\arraytext}[1]{\noalign{\hspace*{-1.75em}
\mbox{#1}\vspace*{-9\jot}}\nonumber\\}
